\clearpage
\section*{Podsumowanie}
\addcontentsline{toc}{section}{Podsumowanie}  

\subsection*{Realizacja celu pracy}
\addcontentsline{toc}{subsection}{Realizacja celu pracy}    

Celem pracy było zaprojektowanie i zaimplementowanie aplikacji sieciowej umożliwiającej analizę biomechaniki martwego ciągu na podstawie nagrania wideo. W ramach pracy przygotowano system składający się z aplikacji klienckiej, aplikacji serwerowej oraz modułu analizy obrazu. Poszczególne części systemu zostały połączone w jeden przepływ działania, w którym użytkownik przesyła nagranie, serwer uruchamia analizę, a wynik zostaje przedstawiony w interfejsie aplikacji.

Zrealizowany system wykorzystuje model estymacji pozy człowieka do wykrywania punktów kluczowych sylwetki. Na podstawie tych punktów obliczane są wybrane parametry ruchu, między innymi kąty związane z ustawieniem biodra, kolana i tułowia. Dane te są następnie wykorzystywane do wykrywania powtórzeń oraz oceny wybranych elementów techniki wykonania ćwiczenia.

W części implementacji wykonano moduł odpowiedzialny za analizę nagrania klatka po klatce, wykrywanie powtórzeń, przygotowanie podsumowania każdego ruchu oraz ocenę techniki na podstawie przyjętych reguł. Zaimplementowano również aplikację serwerową, która przyjmuje plik wideo, zapisuje informacje o analizie i uruchamia moduł przetwarzania. Aplikacja kliencka umożliwia przesłanie nagrania oraz zapoznanie się z wynikiem analizy w formie czytelnej dla użytkownika.

Zakres pracy został zrealizowany zgodnie z jej aplikacyjnym charakterem. Głównym rezultatem nie jest nowy model uczenia maszynowego, lecz działający system informatyczny integrujący gotowy model detekcji punktów kluczowych z własnymi regułami oceny techniki oraz interfejsem użytkownika. Wykonane rozwiązanie pozwala na analizę nagrania martwego ciągu, wykrycie liczby powtórzeń, obliczenie podstawowych parametrów biomechanicznych oraz wskazanie wybranych nieprawidłowości w wykonaniu ruchu.

Przeprowadzone testy potwierdziły, że system realizuje główny scenariusz użycia założony w pracy. W przypadku nagrania przygotowanego zgodnie z wymaganiami, czyli przy widoku z boku, widocznej całej sylwetce i stabilnym kadrze, aplikacja poprawnie wykryła powtórzenia oraz zwróciła ocenę zgodną w dużej części z oceną wizualną. Otrzymane wyniki pokazują, że przygotowany prototyp może pełnić funkcję narzędzia wspierającego wstępną analizę techniki martwego ciągu.

\subsection*{Wnioski z implementacji i testów}
\addcontentsline{toc}{subsection}{Wnioski z implementacji i testów}

Implementacja systemu pokazała, że podział rozwiązania na aplikację kliencką, aplikację serwerową oraz osobny moduł analizy wideo był właściwy dla przyjętego problemu. Taka struktura pozwoliła oddzielić obsługę użytkownika i komunikację z serwerem od właściwego przetwarzania nagrania. Dzięki temu moduł analizy mógł zostać rozwijany i testowany niezależnie od interfejsu aplikacji.

Istotnym elementem pracy okazało się wykorzystanie gotowego modelu estymacji pozy człowieka. Zastosowanie modelu MediaPipe pozwoliło skupić się na interpretacji punktów kluczowych, a nie na samym trenowaniu modelu rozpoznającego sylwetkę. W praktyce oznaczało to możliwość przygotowania własnych reguł służących do wykrywania powtórzeń, obliczania kątów oraz oceny wybranych elementów techniki martwego ciągu. Jednocześnie potwierdziło to, że jakość dalszej analizy jest bezpośrednio zależna od poprawności punktów zwróconych przez model.

Przygotowane reguły oceny techniki sprawdziły się najlepiej w przypadkach, w których nagranie spełniało założenia przyjęte podczas projektowania systemu. Dla materiału wykonanego z boku, przy widocznej całej sylwetce i stabilnym kadrze, system poprawnie wykrył liczbę powtórzeń oraz prawidłowo rozpoznał ruch wykonany w zbyt krótkim zakresie. Wskazane zostały również przyczyny błędu, takie jak obniżony zakres ruchu sztangi oraz brak pełnego wyprostu biodra i kolana.

Testy pokazały również, że system może sygnalizować wybrane problemy z ustawieniem tułowia. W przypadku powtórzenia z widocznym zaokrągleniem pleców aplikacja oznaczyła ruch jako wymagający uwagi. Mechanizm ten nie stanowi pełnej oceny ułożenia kręgosłupa, ale pozwala wychwycić sytuacje, w których relacja barku i biodra w pozycji startowej odbiega od oczekiwanego układu.

Największe ograniczenia ujawniły się w nagraniach odbiegających od założeń systemu. Ustawienie kamery pod kątem utrudniało ocenę ustawienia pleców oraz mogło wpływać na wykrycie pełnej liczby powtórzeń. Z kolei ubiór zasłaniający kolana powodował niestabilną lokalizację punktów kluczowych kończyn dolnych w pojedynczych klatkach. W takich przypadkach ogólna klasyfikacja powtórzenia mogła pozostać poprawna, ale wartości szczegółowych parametrów, szczególnie kątów kolana, wymagały ostrożnej interpretacji.

Na podstawie przeprowadzonych testów można stwierdzić, że system działa poprawnie jako prototyp narzędzia wspierającego analizę techniki martwego ciągu. Najbardziej wiarygodne wyniki uzyskuje się wtedy, gdy użytkownik przygotuje nagranie zgodnie z wymaganiami opisanymi w pracy. Testy nie miały charakteru statystycznej walidacji skuteczności modelu, dlatego otrzymanych wyników nie należy traktować jako pełnej oceny dokładności metody. Stanowią one jednak potwierdzenie, że zaprojektowane rozwiązanie realizuje założony przepływ działania i pozwala na praktyczne wykorzystanie detekcji punktów kluczowych w analizie wybranego ćwiczenia siłowego.

\subsection*{Ograniczenia rozwiązania i możliwości rozwoju}
\addcontentsline{toc}{subsection}{Ograniczenia rozwiązania i możliwości rozwoju}

Przygotowany system spełnia założenia pracy, jednak jego działanie należy rozpatrywać z uwzględnieniem przyjętych ograniczeń. Najważniejsze z nich wynika z wykorzystania analizy obrazu dwuwymiarowego. System ocenia ruch na podstawie punktów kluczowych wykrytych w pojedynczym nagraniu wideo, dlatego nie ma dostępu do pełnej informacji przestrzennej o ustawieniu ciała. W praktyce oznacza to, że wynik analizy jest silnie zależny od ustawienia kamery, widoczności sylwetki oraz jakości detekcji punktów kluczowych.

Ograniczeniem jest również zakres ocenianych elementów techniki. System analizuje wybrane parametry, takie jak zakres ruchu sztangi, wyprost biodra i kolana, pozycję startową oraz uproszczoną relację barku i biodra. Nie stanowi to pełnej oceny biomechanicznej martwego ciągu. W szczególności rozwiązanie nie ocenia wszystkich aspektów techniki, takich jak rozkład obciążenia na stopach, napięcie mięśniowe, rzeczywisty tor sztangi w przestrzeni czy ustawienie kręgosłupa w sposób porównywalny z oceną specjalisty.

Wyniki testów pokazały, że system działa najlepiej w warunkach zgodnych z wymaganiami przyjętymi podczas projektowania. Nagrania wykonane pod niekorzystnym kątem, z częściowo zasłoniętą sylwetką albo z ubiorem utrudniającym rozpoznanie stawów mogą prowadzić do mniej stabilnych wyników. W takich przypadkach błędnie wykryte punkty kluczowe wpływają na wartości kątów oraz komunikaty generowane przez system. Z tego powodu wynik analizy powinien być traktowany jako wskazówka pomocnicza, a nie jako jednoznaczna diagnoza techniki wykonania ćwiczenia.

Kolejnym ograniczeniem jest sposób weryfikacji rozwiązania. Testy przeprowadzono na wybranych nagraniach przygotowanych na potrzeby pracy. Pozwoliło to sprawdzić poprawność działania aplikacji oraz zachowanie systemu w kilku istotnych przypadkach, ale nie stanowi pełnej walidacji statystycznej. Do dokładniejszej oceny skuteczności należałoby przygotować większy zbiór nagrań, obejmujący różne typy sylwetek, ustawienia kamery, warunki oświetleniowe, poziomy zaawansowania ćwiczących oraz przykłady różnych błędów technicznych.

Możliwości dalszego rozwoju systemu obejmują przede wszystkim rozszerzenie i dokładniejszą walidację reguł oceny techniki. Warto byłoby porównać wyniki działania aplikacji z oceną trenera lub fizjoterapeuty oraz na tej podstawie dostosować progi wykorzystywane w klasyfikacji powtórzeń. Istotnym kierunkiem rozwoju byłoby także poprawienie odporności systemu na błędne wykrycia pojedynczych punktów kluczowych, na przykład przez dodatkowe wygładzanie danych, analizę pewności detekcji oraz odrzucanie klatek o niskiej jakości.

Rozwiązanie można również rozbudować od strony aplikacyjnej. Przydatne byłoby dodanie dokładniejszych instrukcji przygotowania nagrania, podglądu poprawności kadru przed analizą oraz bardziej szczegółowej prezentacji przebiegu ruchu w czasie. W dalszej perspektywie możliwe byłoby także rozszerzenie systemu o analizę innych ćwiczeń siłowych lub wykorzystanie bardziej zaawansowanych metod analizy ruchu, na przykład nagrań z kilku kamer albo modeli zwracających dokładniejsze informacje przestrzenne.

Podsumowując, wykonany system stanowi działający prototyp aplikacji wspierającej analizę techniki martwego ciągu na podstawie nagrania wideo. Praca pokazała, że połączenie gotowego modelu estymacji pozy człowieka z własnymi regułami oceny może być użyteczne w prostych warunkach treningowych. Jednocześnie przeprowadzone testy wskazały, że wiarygodność takiej analizy zależy od jakości nagrania oraz poprawności wykrycia punktów kluczowych. Dalszy rozwój rozwiązania powinien więc koncentrować się zarówno na rozszerzeniu walidacji, jak i na zwiększeniu odporności systemu na trudniejsze warunki nagrywania.